% §13.3 The Composition Functor
%
% Monograph register + the sum-exception fix. Bernoulli is a structure-preserving
% map that commutes with the product, list, and function constructors (Kronecker);
% the SUM is the exception (tag coupling, does not factor). Adds rem:monad (the
% Kleisli structure salvaged from ex:ch13-6). Cross-refs thm:bernoulli-product,
% prop:sum-nonfactor, ch:set-composition, ch:sketches-estimators,
% ch:approximate-relations, ch:approximate-functions. DEFINED:
% sec:composition-functor, rem:monad.

\section{The Composition Functor}
\label{sec:composition-functor}

\Cref{sec:bernoulli-of-composite} gave the lift on each constructor. This section
observes that the rules belong together as one structure-preserving map, with the
sum as the documented exception.

\subsection*{A structure-preserving map}

The lift $\Bernoulli{-}$ takes a type-level operation to a channel-level one. The
product becomes a Kronecker product (\Cref{thm:bernoulli-product}), the list a
Kronecker power, the function a per-input Kronecker product. A map that takes
operations to corresponding operations is a \emph{structure-preserving map}; the
formal categorical name is \emph{functor} (\S13.5). The four cases:
\begin{center}
\begin{tabular}{lll}
\hline
Type level & Channel level & Reading \\
\hline
$T \times U$ & $\Bernoulli{T} \otimes \Bernoulli{U}$ & Kronecker product \\
$L T$ & $L \Bernoulli{T}$ (mixture if length noisy) & per-position lift \\
$T \to U$ & $T \to \Bernoulli{U}$ & per-input lift \\
$T + U$ & does \emph{not} factor & tag coupling (\Cref{prop:sum-nonfactor}) \\
\hline
\end{tabular}
\end{center}
The first three rows are isomorphisms: the type and the channel specify the same
$\Bernoulli{-}$ instance up to relabeling, and $\Bernoulli{-}$ commutes with the
product, list, and function constructors. The fourth row is the exception, the
sum, where the tag error couples the components and no factorization recovers the
channel. That exception is not a gap in the theory but its foundation: the
optional type's tag error is the atomic membership rate (\S13.2).

\begin{remark}[$\Bernoulli{-}$ as a monad]
\label{rem:monad}
The function row $\Bernoulli{T \to U} \cong T \to \Bernoulli{U}$ is the signature
of a \emph{monad}. The Kleisli arrow, a noise-introducing function $T \to
\Bernoulli{U}$, takes a pure value to a noisy one; two arrows $f : T \to
\Bernoulli{U}$ and $g : U \to \Bernoulli{V}$ compose into $T \to \Bernoulli{V}$
by feeding $f$'s output through $g$ and averaging the intermediate noise (the
monad's bind); and a pure value embeds as the point distribution concentrated on
it (the unit). These three make $\Bernoulli{-}$ a monad in the
programming-language sense, the noisy-value counterpart of the optional and list
monads, with the formal axioms in \textcite{wadler-1992}.
\end{remark}

\subsection*{Why it matters}

Without the table, every approximate structure needed its own composition
derivation: chapter~6 for sets, chapter~10 for sketches, chapter~12 for
relations. With it, composition is mechanical for any type built from the three
factoring constructors: decompose the type and apply the matching row at each
level. Adding a new structure becomes a matter of identifying its type, not
redoing the algebra, a sketch estimating matrix entries is
$\Bernoulli{T \times U \to \mathbb{R}}$, and the function and product rows
determine its composition with no fresh derivation.

The structures already built each instantiate the table. An approximate set
(\Cref{ch:set-composition}) is $\Bernoulli{T \to \BoolType}$, so the function row
gives the composition rule and chapter-6's results are the function row at
codomain $\BoolType$, read pointwise. A CountMin sketch
(\Cref{ch:sketches-estimators}) is $\Bernoulli{T \to \mathbb{N}}$, its additivity
the product-row Kronecker structure on the codomain. An approximate relation
(\Cref{ch:approximate-relations}) is $\Bernoulli{T \times U \to \BoolType}$,
combining the function row (per-input lift) and the product row (on the domain).
The n-ary channel of \Cref{ch:approximate-functions} is the building block $X \to
Y$ underneath them all, distinguished only by codomain and domain structure.

\subsection*{The framework as one thing}

There is one underlying object across the book. The atom is the noisy bit
$\Bernoulli{\BoolType}$; every other structure is a $\Bernoulli{T}$ for an
algebraic type $T$ built from $\BoolType$ and the constructors, and the
composition rules that organized the earlier chapters are the rows of the table.
The title's ``algebra of approximate computation'' is the algebra of the type
constructors, with $\Bernoulli{-}$ its action on the approximate side, factoring
through the product, list, and function and meeting its irreducible limit at the
sum's tag error. The remaining chapters specialize the algebra rather than extend
it: chapter~14 builds approximate sets at the entropy floor, chapter~15 the
perfect-hash filters, chapter~16 catalogues where the independence assumptions
break, and chapters~17 and~18 turn the algebra to cryptographic use. Each is a
specialized application of the table.
